\chapter{The exhaustive canonical source cover}
\label{ch:cover}
\label{component:pe1088004af65}
% Fresh exhaustive canonical-source assembly, 2026-10-22.

\label{pe1088004af65:sec:complete-source-routing}

\paragraph{Quantifiers and the full-parent target.}
Let \(K_c,K_t\) be arbitrarily chosen, separately labelled real closed
extensions of \(\mathbb R\).  There is no algebraic-disjointness
assumption: the two factors may be labelled copies of the same field.
Put \(T=K_c\otimes_{\mathbb R}K_t\).
Consider a physical canonical source over \(K_c\) with an actual
finite noncommuting stationary zero and one of the exact patterns
\[
\begin{array}{ll}
{\rm C1}:&d,m,C,A,\alpha>0,\quad\lambda=n=0,\quad
             D,\beta\geq0,\quad D+\beta>0,\\
{\rm C2}:&d,m,C,A,D,\alpha>0,\quad\lambda=\beta=0,\quad n\geq0.
\end{array}\tag{LC1}
\]
These are precisely the two elementary source hypotheses produced
by the extreme-source lemma and SG10 in Chapter~\ref{ch:geometry}.  For finite independent
testing half-angles and any state \((W,X,Y,Z)\in K_t^4\), the complete
parent is
\[
\begin{split}
 F={}&4\lambda D_u(W+vX)^2+4nD_uD_vX^2
       +4m(vW-X+uvY-uZ)^2\\
 &+4CD_v(uW-Y)^2+4AD_uD_vY^2
       +(4D+2d)(uW+uvX-Y-vZ)^2\\
 &+(4\alpha+2d)D_u(vY-Z)^2
       +(4\beta+2d)D_v(X+uZ)^2-2dD_uD_vZ^2 .
\end{split}\tag{LC2}
\]
Every coefficient times a testing expression denotes its separated
tensor product.  The homogeneous angle form of LC2, rather than
an affine limit obtained by a cancelled zero coordinate, defines
the literal projective endpoints.

\begin{theorem}[Complete local certificate]
For every source in LC1 and every independent projective testing pair
and real or complex testing state, the entire homogeneous parent is
a finite sum of squares in \(T\).
This includes every coefficient boundary permitted in LC1, every
positive source opening, all relative source and testing rates, and
all state ratios.  The source contact is never identified with the
testing point.
\end{theorem}

\section{Normalization and actual native entry}

Divide by the positive pure-source \(d\).  This operation is reversible
in a tensor sum of squares by multiplication of each root by
\(\sqrt d\otimes1\).  Before or after this normalization the invariant
total-scale quantity is
\[
 \delta_{\rm tot}
  =\frac d{d+\lambda+n+m+C+A+D+\alpha+\beta}.
\tag{LC3}
\]
It is positive and at most one.  Its positive-unit and positive
infinitesimal alternatives are exhaustive.  In the first alternative
the complete positive-scale theorem AU1--AU19 applies after EN.
The second alternative retains the hypothesis
\(\delta_{\rm tot}\in\mathrm{Inf}_{c,>0}\) in every residual theorem
used below.  Setting \(d=1\) does not erase that hypothesis.

The theorem EN1--EN16 derives the full native signs and exact rational
source identification from physicality and the actual finite contact.
Its proof uses the contact lemma and SG17--SG22 in Chapter~\ref{ch:geometry}: coercive Laurent support, the strict single-field contact Hessian, contact sufficiency, positive $W_c$ after $Z_c=1$ normalization, and the event-row deletion implication.
It does not assume any tensor certificate being assembled here.
For C1, after the exact Alice source reflection, its native coordinates
satisfy
\[
\begin{gathered}
 0<k\leq1,\quad \eta,\zeta>0,\quad
 r=1-a\eta\in(0,1),\quad \theta=t\zeta\in(0,1),\\
 0<a,t<2,\quad x=k^2,\quad h=1-\theta,\quad
 E=1-x\eta h>0,\quad Q=1-x\zeta r>0,\\
 M=1+xr-t>0,\quad J=1+xh-a>0,\qquad
 e=k\sqrt{\eta/\zeta}>0 .
\end{gathered}\tag{LC4}
\]
The original \(D,\beta\) remain nonnegative, and either may be zero.
For C2, EN gives exactly CT2 and its unique actual stationary
coefficient formulas, including \(n=0\).
All these implications hold over every \(K_c\) by the first-order
transfer proved in EN; no topological compactness of a
non-Archimedean field is being presumed.

\section{Every weak-total C1 source}

Use the native parameters LC4, the actual original coefficients, and
\(\delta_{\rm tot}\) positive infinitesimal.
Here I, U and O mean a positive infinitesimal, positive bounded unit,
and positive unbounded source quantity, respectively.
Since \(0<k\leq1\), \(k\) has type I or U.
The following disjoint scale dispatch is exhaustive:
\[
\begin{array}{c|c|c|l}
k&\zeta&\text{remaining source type}&\text{actual theorem}\\ \hline
{\rm U}&\text{any}&\eta,\zeta:{\rm I}&{\rm AS}\\
{\rm U}&\text{any}&\eta,\zeta\text{ not both I}&
                      \text{ordinary-Alice join below}\\
{\rm I}&{\rm U\ or\ O}&\eta>0&
                      {\rm JC}\text{ with bare source scope}\\
{\rm I}&{\rm I}&\eta:{\rm I}&
                      \text{uniform all-}e\text{ join}\\
{\rm I}&{\rm I}&\eta:{\rm O}&{\rm C1.1}\\
{\rm I}&{\rm I}&\eta:{\rm U},\ e:{\rm I}&{\rm C1.3}\\
{\rm I}&{\rm I}&\eta,e,r:{\rm U}&{\rm BE}\\
{\rm I}&{\rm I}&\eta,e:{\rm U},\ r:{\rm I}&{\rm CA}\\
{\rm I}&{\rm I}&\eta,r:{\rm U},\ e:{\rm O}&{\rm UE}_{\rm old}\\
{\rm I}&{\rm I}&\eta:{\rm U},\ e:{\rm O},\ r:{\rm I}
                                                  &{\rm C1.2}.
\end{array}\tag{LC5}
\]
The rows using \(r\) are exhaustive because \(0<r<1\).
The positive number \(e\) has exactly the three types shown.
No comparison with an independent testing quantity occurs in LC5.

\paragraph{Ordinary Alice: the actual UI, UU and UH split.}
Suppose \(k\) is a unit and \(\eta,\zeta\) are not both infinitesimal.
If \(\eta\) is infinitesimal, then \(r\sim1\) and \(Q>0\)
make \(\zeta\) bounded.  Thus \(\zeta\) is a unit in this case.
If \(\zeta\) is infinitesimal, then \(\theta\) is infinitesimal,
\(h\sim1\), and \(E>0\) make \(\eta\) bounded; hence \(\eta\)
is a unit.  In particular the two possibilities
\[
 (\eta:{\rm I},\zeta:{\rm O}),\qquad
 (\eta:{\rm O},\zeta:{\rm I})
\tag{LC6}
\]
are impossible when \(k\) is a unit.
Since \(v_c^2=x\eta\zeta\), the remaining possibilities are:
one of \(\eta,\zeta\) infinitesimal and the other a unit, giving
actual UI; both units, giving actual UU; or both noninfinitesimal
with at least one unbounded, giving actual UH.

The actual UI source theorem and its exact native stabilizer cover
the first possibility.  More explicitly, the exact AS3--AS5
full-parent reflection reduces it to \(\eta\) a unit and
\(\zeta\) infinitesimal.  This reflection preserves the C1 omission
pair and the original nonnegative coefficients; its new native
parameters and gaps are checked in that identity.  Its Alice
parameter remains a unit, but a second assertion \(k\leq1\) is
not imposed on the reflected representative.
Use the complete CU1--CU13 theorem in
Section~\ref{component:pa0b854a04f49}, together with CR1--CR13 in the actual
complementary-testing body.  For these displayed CU source
quotients rescale by the positive pure-source factor to its
normalization \(C=2\); the native parameters are unchanged,
and this local normalization is restored afterwards.
In CU put \(g^2=\zeta\).
The exact source quotients make \(C/a,A/a\) strong,
\(d/a,\alpha/(ag^2)\) units, and \(m/a,\beta/a\) large.
CU2 retains the entire paired \(m,\beta\) normal packets;
CU3 retains the centered \(C,A,D\) excesses.
CU5 is the whole remaining quadratic state polynomial on
its bounded six-root frame.  CU6--CU9 explicitly divide all
complete first-normal coefficients and the scalar correction in the
contact ideals.  Their exact tensor zero redistributes the
two angular pivots; CU10 computes the full matched state
Schur Hessian.  Its positive determinant uses the original
\(D\geq0\) numerator.  The possible divergent \(A,D\) pair has
independent hard rows, so its constrained finite-cap limit
keeps the angular components of those rows.
At unmatched Alice residue the normalized Laurent coefficients
retain coercivity and strict convexity.  The actual real
two-block lemma UA1--UA3 is applied using the exact relation
\(L(gZ)-Q(vZ)=0\) between its root products, not by inverting
a determinant in \(T\).  CR supplies all complementary tests,
especially its CR5a five-root reciprocal-Bob frame retaining
the weak \(\alpha\) packet.  This proves all UI tests and states.
Literal \(D=0\) is retained; \(\beta=0\) is impossible here by
the displayed lower bound on the actual \(\beta/a\).
For actual UU use the full NC1--NC7 source proof in
Section~\ref{component:p42ee07a067c1}.
It first proves the uniform native gap \(E\geq1/18\).
The nonnegative \(D,\beta\) formulas and their exact elimination
identity give
\[
 C,A,m,\alpha\text{ strong},\qquad
                    D\text{ or }\beta\text{ strong}.
\tag{LC7}
\]
In more detail \(x,v_c^2,\eta,\zeta\) are units; the specialization
\(B_\beta|_{\theta=1}=-1-x\zeta(1-r)\) makes \(h\) a unit.
The original numerator polynomials \(P_4,P_\alpha,P_d\) are
positive units.  If both \(D,\beta\) were weak, the left side of
the exact NC7 elimination would be infinitesimal and its
explicit positive right side a unit.  This is impossible.
The five consecutive strong events in LC7 therefore give the
actual five-event all-tests proof, with its detailed A6
noncommuting testing complement.
This is a source theorem at all upper event scales, not an
application of a bounded-coefficient theorem to unbounded data.
For actual UH the native NI8--NI10 implications give two
alternatives: \(\eta\) unbounded with \(\zeta\) a unit, or
\(\zeta\) unbounded with \(\eta\) noninfinitesimal.
The exact \(\beta\)-numerator identity proves \(\eta h\) a
unit in both.  Thus, with \(\sigma=1/\eta\), \(z=1/\zeta\)
and \(H=h/\sigma\), the source is respectively
\(\sigma\) infinitesimal and \(z\) a unit, or \(z\)
infinitesimal and \(\sigma\) bounded, always with \(H\)
a unit.  These are precisely the printed UH1 alternatives;
no bound on the ratio of two source infinitesimals is added.
The actual four-consecutive-event theorem CE and its
matched weak outer frames CE17--CE21 remove the remote tests.
The remaining strip has unit testing \(u\) and infinitesimal
\(w=1/v\).  The complete homogeneous biquadratic conductor
Section~\ref{component:pf7ffa921793c} applies there.
With \(b=r/z\), use its scaled weights
\(\widehat m=m/z\), \(\widehat C=C/z\),
\(\widehat A=Az\) and \(\widehat\alpha=\alpha z\).
Their products
\(\widehat mM,\widehat C(1-r),\widehat A b,\widehat\alpha Q\)
are bounded and cancel the opening poles before the full
scalar and first-normal collection.
At a matched Alice residue its exact angular zero and
the full hard Schur calculation redistribute the positive
angular pivots.  At an unmatched residue its raw angular
product relation gives the finite real two-block allocation.
The single and double opening limits retain all divergent
centered rows; their independent hard-row determinants and
bounded hard minimizers justify the constrained limits.
This is a full-parent construction, including the
\(J=\theta=0\) residue face.
Both \(D=0\) and \(\beta=0\) are retained where permitted by the
actual source equations.  Thus no assertion that UI or UH alone
already covers UU is needed.

\paragraph{The uniform source and the four earlier exits.}
In the fourth row of LC5, \(k,\sqrt\eta,\sqrt\zeta\) are all
infinitesimal, while \(e=k\sqrt{\eta/\zeta}\) is arbitrary.
The new uniform all-\(e\) join UE1--UE4 applies exactly.
For \(e\) infinitesimal it uses the fresh complete UJ proof,
with its full bounded-ratio and inverse-ratio state charts.
Otherwise \(1/e\) is bounded, and the actual complete
bounded-Delta proof supplies its restricted monic module,
weighted full scalar, both entire first normals, hard block,
and complementary state chart.  Its outer current formulas
are joined in UE without an unprinted restriction on \(e\).
Neither a historical whole-uniform assertion nor a missing
old conductor label is a premise.

The AS, BE, CA and \({\rm UE}_{\rm old}\) rows use their actual
full-parent proofs, respectively
Section~\ref{component:p9f00f949fc8c},
Section~\ref{component:p69995d829bb2},
Section~\ref{component:p3763434080c2} and
Section~\ref{component:p890eea5f538b}.
Their ordinary-angle, centered-A, contact-conductor and
outer-current subordinate bodies are included explicitly.
AS proves the needed strict original event signs on its source
domain.  BE and CA include \(\beta=0\).
The last proof derives \(\beta/d\) unbounded, so its intersection
with literal \(\beta=0\) is empty by an actual source estimate.
These proofs use the bare LC4 equations and their stated source
scales; no estimate comes from exclusion of an unnamed earlier
admission.

\paragraph{The remaining C1 rows and the complete C1.4 union.}
For C1.1 the actual source reduction, including the printed
\(k^2\eta\leq1/2+\mathrm{Inf}_c\) derivation, splits into
the unit-\(q\) and infinitesimal-\(q\) branches.  The complete
unit-\(q\) body and the protected-normal local and remote
infinitesimal-\(q\) bodies prove every testing/state chart.
C1.2 uses the actual centered-A bounded-Delta theorem, including
its full contact conductor and centered-A outer proof.
C1.3 uses the actual small-\(e\) all-tests theorem and its
explicit outer-deficit calculation.  These are full bare-source
statements with the scales printed in LC5.

For the third row of LC5, the source is exactly the native
C1.4 input \(k\) infinitesimal, \(\zeta\) noninfinitesimal,
with all original gaps and \(\delta_{\rm tot}\) infinitesimal.
We now give its bare-source join explicitly; this also explains
why a blanket earlier-complement predicate in the entered JC
statement is unnecessary for this invocation.

Put \(L=x\zeta\).  If \(L\) is unbounded, \(Q=1-Lr>0\)
forces \(r\) infinitesimal.  Since \(1-r=a\eta\) and
\(0<a<2\), \(\eta\) cannot be infinitesimal.
Unit \(\eta\) is the complete Section~\ref{component:pd4b92bf7de71}
exit.  For unbounded \(\eta\), put
\(\sigma=1/\eta\), \(H=h/\sigma\) and
\(M_\beta=\beta\theta/\sigma^2\).
Bounded \(H,M_\beta\), including \(\beta=0\), is the actual
JR/JB theorem.  Bounded \(H\) with unbounded \(M_\beta\)
is the full C14.1a SF conductor.  Unbounded \(H\) is
the C14.1b source: SF first treats \(\operatorname{res}h<1/2\);
US bounds \(k/\sigma\) and treats its unit case.
The BS/BC/GA and CE testing and state exits then leave
exactly the small-sphere source and the infinitesimal
testing chart stated in JC5.  The fresh AR theorem restores
the angular budget and all SP6 first normals there.
AR requires no bound on \(qu/w\), \(u/w\), \(W/Z\) or
\(uW/(wZ)\).  This closes the large-\(L\) alternative
without a comparison of independent source and test rates.

If \(L\) is a unit, the three types of \(\eta\) give,
respectively, the complete Section~\ref{component:p080aa62de16d}
theorem, the LU theorem in Section~\ref{component:p686209b87229},
and C14.2.  The first two use the actual LII parent
dictionary with their own native coefficient checks.
They do not identify the source with an arbitrary C2
source satisfying a narrower theorem.
C14.2 uses its explicit NC/LC source bounds,
complete \(C\)-opening extraction, HG/CT radial matrix,
sphere-weighted budgets, CS/CF/HG state joins and LA/LB
remote tests.  Every part has the full parent and all
normal terms.

If \(L\) is infinitesimal and \(\eta\) unbounded, then
\(\sigma=1/\eta\) is infinitesimal and this is C14.4.
Its source-only parameters and exact identity are
\[
\begin{gathered}
 S=\sigma+L,\quad \kappa=x/S,\quad
 \nu=\sigma/S,\quad \ell=L/S,\quad \nu+\ell=1,\quad
 \zeta=\ell/\kappa,\\
 Q_s=1-\ell Sr,\qquad
 \kappa Q_s+\ell M+\nu E=1,\qquad E\geq1/18.
\end{gathered}\tag{LC4a}
\]
Consequently \(\kappa\) is bounded.  The fractions
\(\nu,\ell\) cannot both be infinitesimal; if \(\kappa\)
is a unit, then \(\ell=\kappa\zeta\) is a unit.
Thus the five exact source cells are
\[
\begin{array}{c|c|c|l}
\kappa&\nu&\ell&\text{whole-parent theorem}\\ \hline
{\rm U}&{\rm U}&{\rm U}&{\rm FU}\text{ with its UKR matched chart}\\
{\rm U}&{\rm I}&{\rm U}&{\rm NU}\\
{\rm I}&{\rm U}&{\rm U}&{\rm FM}\\
{\rm I}&{\rm U}&{\rm I}&{\rm EL}\\
{\rm I}&{\rm I}&{\rm U}&{\rm BN,\ UN}.
\end{array}\tag{LC4b}
\]
For the last cell put \(q_\nu=\nu/\ell\),
\(\tau=h/q_\nu\) and \(n_s=\beta/q_\nu^2\).
Bounded \(\tau\) gives BN for every allowed \(n_s\);
unbounded \(\tau\) gives UN, including infinitesimal
\(\theta\). The shared SN coefficient dictionary is retained,
but a separate bounded-\(n_s\) terminal is unnecessary. Each source is covered.
All these proofs start from the exact FM dictionary,
NC5 and their actual source/sphere equations.
Their SE/EB/CE testing exits have those same printed
source inequalities.  They retain all normal terms,
literal coefficient zeros where permitted, and every
source opening.

It remains to take \(L\) infinitesimal and \(\eta\)
bounded.  Unbounded \(\zeta\) is C14.3a.
Unit \(\zeta\) with infinitesimal \(\eta\) is C14.3b.
Their complete JF parent is combined with DG/LD normal
conductors and the completed \(J\)-row conductor;
the latter removes the earlier \(x/\eta\) restriction.
The old restricted follow-up assertion is not invoked.
The completed-J conductor, the pure-testing DS3 state exit and SB supply all states and the remote tests.
The small-\(\eta\), bounded-\(a\) input is used only
through its proved angular paragraph, not its old
proposed radial continuation.

The remaining \(\eta,\zeta\) are units.  Use the
actual positive source gap
\[
 j=\eta(1+xh)-1+r=\eta J>0,\qquad
 j_0=\eta+r-1,\qquad j-j_0=\eta xh\in\mathrm{Inf}_c .
\tag{LC4c}
\]
It is bounded.  Unit \(j\) gives positive unit \(j_0\)
and the RO theorem.  Infinitesimal \(j\) gives
\(\eta+r\sim1\) and unit \(1-r\); unit \(r\) then gives
JO and infinitesimal \(r\) gives HO.
RO uses the IC parent dictionary with its newly proved
native bounds, not the narrower IC1 source premise.
JO and HO use \(j>0\), including \(j_0=0\) or slightly
negative.  Their exact three-circuit body and SB remote
proof retain all normal terms.

This proves the entire \(L,\eta,\zeta,j\) source
partition directly from LC4.  Every local source
estimate just invoked uses these native equations,
the original \(D,\beta\geq0\), its printed scale cell,
the NC gap and, where required, the unchanged
infinitesimal total scale.  No bound is obtained from
failure of an unnamed earlier theorem.
The seven restored exits are exactly large-\(L\)/unit-\(\eta\),
unit-\(L\)/small-\(\eta\), unit-\(L\)/unit-\(\eta\),
JR/JB, RO, JO and HO.  They are positive parts of this
cover.  Thus the actual JC constructions prove the
bare third row of LC5; any original earlier-complement
predicate can afterwards be intersected with the result.
The strengthening preserves physicality, signs, gaps,
literal boundaries and total scale.

All ten rows of LC5 now have complete original-parent
certificates.  They include every residual C1 source after
the positive-total-scale alternative, so C1 in LC1 is proved.

\section{Every C2 source and the absence of a circular premise}

The theorem CT1--CT16 proves the entire native C2 family.
Its proof takes the same total-scale split, treats both
weak complementary events by actual radial source entry,
and otherwise gives all nine ordered actual-angle types
in its explicit table.  Every empty cell has an exact
necessary-source contradiction; every occupied cell has
its complete local and remote full-parent proof.
In particular its directed weak HH cell is reflected to
the exact HJ domain, with \(C\) strong, \(\alpha\) weak
and both actual angles infinitesimal.
The HJ proof includes all ten source rows and every
allowed \(n=0\) map.  Its C1 boundary calls use only the
two explicit centered-A and small-\(e\) bare-source bodies,
never the C1 conclusion of the present theorem.
Thus CT has no dependence on this global C1/C2 join.
This proves C2 in LC1.

The mathematical dependencies have the following direction:
\[
\begin{gathered}
 \text{actual physical source geometry}
       \ \Longrightarrow\ \text{the source/contact reduction and EN},\\
 \text{local identities and strict tensor Gram}
       \ \Longrightarrow\ {\rm AU,\ JC,\ HJ,\ UJ},
       \ \text{scoped terminals},\\
 \text{these terminals and finite source partitions}
       \ \Longrightarrow\ {\rm CT,\ LC}.
\end{gathered}\tag{LC8}
\]
A6 is separately proved from the actual zero-\(d\) geometry,
one- and two-hard-packet cases, weighted matching, promotion
and induction bodies; it does not use LC, CT, JC or HJ.
EN does not use AU or any tensor-SOS conclusion.
The source entries and coefficient-field physicality arguments
inside local proofs are not converted into positivity
statements in the tensor ring.
At no step is a mixed Schur denominator cancelled.
All roots and all first and higher state-normal terms belong
to the actual full parent LC2.

The finite mathematical inventory accompanying this proof
records direct and transitive actual proof bodies with their
restricted imported statements.  An old open-complement
paragraph, a historical status, a filename, or a checker
assertion count is not a mathematical premise.
Finite typed-source checks corroborate the exhaustion of
LC5 and CT6; strictness, conductor membership and full-normal
restoration are proved in the displayed mathematical bodies.

\section{Restoration and the universal slack consequence}

Every local proof gives the unmultiplied parent, or an
equivalent parent multiplied only by a positive pure-source
or pure-testing factor whose nonvanishing is explicitly
proved.  Such a factor is removed by its square root inside
that field.  Original retained packets are then added with
their actual nonnegative coefficients.
The exact homogeneous angle identities restore every finite
chart and include literal endpoints without division by
a vanishing coordinate.  Pure-testing state normalizations
cover all nonzero real projective states, and the zero state
is immediate.  Since LC2 is a real quadratic form, applying
the real result to real and imaginary components and adding
proves the complex-state assertion.  Restoring \(d\)
finishes the complete local-certificate theorem.

\begin{corollary}[Universal separated slack]
Let \(\mathcal Q_{222}\subset\mathbb R^8\) be the finite-dimensional
bipartite tensor-product quantum set in the two-setting,
two-outcome scenario, and let \(P\) be its cone of nonnegative
affine functionals.  Use the fixed finite real-qubit formula
for \((\mathcal Q_{222})_{K_t}\) and the transferred cone
\(P_{K_c}\).  Then for every separately labelled \(K_c,K_t\)
as above,
\[
 f_0\otimes1+\sum_{\nu=1}^8 f_\nu\otimes q_\nu
       \ \in\ \Sigma T^2
 \qquad
 (f\in P_{K_c},\ q\in(\mathcal Q_{222})_{K_t}).
\tag{LC9}
\]
\end{corollary}

\begin{proof}
The complete source reduction in Chapter~\ref{ch:geometry} writes every normalized
source in \(K_c\) as a finite conic combination of event atoms,
the two explicit noncoercive closed faces, and the actual
canonical contact sources LC1.  It proves this fixed finite
statement over the reals and transfers it; no arbitrary
infinite extreme-point expansion is used.
Event atoms have separated squares.  Each of the two closed
faces has the actual finite matrix lift and strict feasible
point in Chapter~\ref{ch:geometry}, equations SG11--SG14.
Attained finite-dimensional duality and separate-field
positive-semidefinite factorizations give its tensor
trace-pairing squares.  Canonical atoms use the theorem
just proved, followed by the exact homogeneous slack
normalization in the physical reduction.

Each testing point is a convex combination of at most nine
real-qubit atoms in its own field.  If the source combination
weights are \(a_j\geq0\), testing weights are \(w_i\geq0\)
with \(\sum_iw_i=1\), and each atom pairing is
\(\sum_k h_{ijk}^2\), then the full pairing is exactly
\[
       \sum_{i,j,k}
        \left[(\sqrt{a_j}\otimes\sqrt{w_i})h_{ijk}\right]^2 .
\tag{LC10}
\]
Affineness preserves the constant term because the testing
weights sum to one.  Zero weights are never inverted.
All Bell relabelings and all positive normalizing factors
are restored inside their own field.  This is finite conic
and convex closure in the tensor ring, and proves LC9.
\end{proof}

The proof supplies the universal tensor statement, not an
evaluated affine matrix pencil or a numerical bound on its
size.  The separate finite-lift criterion applies to LC9
using two labelled copies of any one real closed extension;
its compactness conclusion is not inferred from a sum of
the local root counts.

\section{\texorpdfstring{Every canonical C2 source: an exhaustive mathematical join}{Every canonical C2 source: an exhaustive mathematical join}}
\label{component:pe351eb06b28b}
% Fresh source partition and proof-use audit, 2026-10-22.

\label{pe351eb06b28b:sec:c2-complete-source-routing}

\paragraph{Scope and notation.}
Let \(K_c,K_t\) be arbitrarily chosen, separately labelled real closed
extensions of \(\mathbb R\); no algebraic-disjointness assumption is made.
All source inequalities below are in \(K_c\).
Use the entire original physical C2 parent
\[
\begin{split}
 F={}&H+4nD_uD_vX^2+4mM^2+4CD_v(uW-Y)^2
       +4AD_uD_vY^2+4DJ^2+4\alpha D_u(vY-Z)^2,\\
 M={}&vW-X+uvY-uZ,\qquad J=uW+uvX-Y-vZ,\\
 H/2={}&J^2+D_u(vY-Z)^2+D_v(X+uZ)^2-D_uD_vZ^2 .
\end{split}\tag{CT1}
\]
Here \(D_u=1+u^2,D_v=1+v^2\), \(d=1\), and the omitted weights are
\(\lambda_{\rm event}=\beta=0\).  Assume an actual finite noncommuting
stationary zero.  The native-entry theorem EN gives exactly
\[
 m,C,A,D,\alpha>0,\quad n\geq0,\quad
 c,a,b,t>0,\quad 0<r<1,\quad
 h=b-a^2t>0,\quad g>0,\quad j>0,
\]
\[
 g=(1+c^2)b-(1-r+a^2c^2t),\qquad j=1+c^2r-t,\qquad
 1-c^2h>0,\quad b-a^2c^2r>0.                    \tag{CT2}
\]
The actual angles are \(u_c=c,v_c=ac\), and the six coefficients
are the actual rational stationary coefficients, with the nonzero
stationary determinant proved in EN.  In particular no limiting or
formally perturbed source is substituted for CT1.

\paragraph{Theorem.}
Under CT1--CT2, the full parent is a finite sum of squares in
\(K_c\otimes_{\mathbb R}K_t\) at every projective testing pair and
every real or complex state over \(K_t\).
All coefficient scales, literal \(n=0\), and every positive source
opening allowed by CT2 are included.  This theorem has no
earlier-admission-complement hypothesis.

\subsubsection*{Two finite partitions}

First put
\[
 \delta_{\rm tot}=(1+n+m+C+A+D+\alpha)^{-1}.
\tag{CT3}
\]
If \(\delta_{\rm tot}\) is a unit, all six coefficients are bounded,
and the complete positive-scale theorem AU applies at its now-proved
native-entry interface.  Otherwise \(\delta_{\rm tot}\) is positive
infinitesimal.  No normalization of an individual large event is
used to replace this dichotomy.

Call a positive event weak when it is infinitesimal after this
normalization \(d=1\), and strong otherwise.  A strong event may be
unbounded.  If exactly one of \(C,\alpha\) is weak, orient it to be
\(C\), using the exact full-parent reflection
\[
\begin{split}
 &(u,v;W,X,Y,Z;n,m,C,A,D,\alpha)\\
 &\quad\longmapsto
 (1/v,1/u;Z,X,Y,W;n,D,\alpha,A,m,C).
\end{split}\tag{CT4}
\]
In homogeneous angle coordinates the identity is
\(F=u^2v^2F'\); its polynomial homogeneous interpretation includes
literal testing zero and infinity.  CT4 preserves all source gaps:
writing \(z=c^2\), its parameters are
\[
 c'=(ac)^{-1},\quad a'=a,\quad b'=a^2/b,\quad
 r'=a^2zr/b,\quad t'=(1-zh)/b,
\]
\[
 h'=a^2zh/b,\quad g'=a^2j/b,\quad j'=g/b,\qquad
 \Delta'=\Delta/(b^2z^2),\quad N'=a^2N/b^2 .
\tag{CT5}
\]
These are pure-source identities, and give positive values for
\(1-r',1-z'h',b'-a^2z'r'\).  They preserve
\(\delta_{\rm tot}\) exactly.  EN and HJ verify the full reflection;
the retained matching-source entry prints the displayed rational
source identities.  No assertion \(t\mapsto1/t\) is used.

Let \({\rm I},{\rm U},{\rm O}\) denote, respectively, a positive
infinitesimal, a positive bounded unit, and a positive unbounded
source number.  Every positive actual angle has exactly one of
these types.  After the two-weak case is treated below, the
remaining dispatch is
\[
\begin{array}{c|c|c}
(u_c,v_c)&C,\alpha\text{ both strong}
                 &C\text{ weak},\ \alpha\text{ strong}\\ \hline
{\rm II}&\text{strong II local and remote}&\text{impossible}\\
{\rm IU}&\text{single-angle join}&\text{impossible}\\
{\rm IO}&\text{strong IH local and remote}&\text{weak IH full join}\\
{\rm UI}&\text{strong UI full join}&\text{impossible}\\
{\rm UU}&\text{five retained events}&\text{impossible}\\
{\rm UO}&\text{single-angle join by CT4}&\text{single-angle join by CT4}\\
{\rm OI}&\text{strong HI full join}&\text{impossible}\\
{\rm OU}&\text{strong UI full join by CT4}&\text{impossible}\\
{\rm OO}&\text{strong II full join by CT4}&\text{HJ by CT4}
\end{array}\tag{CT6}
\]
The letters H in the retained file titles mean the type O in this
table.  They denote an unbounded \emph{finite} actual half-angle,
not a literal projective endpoint.  Testing endpoints are included
by each full-parent theorem.

\subsubsection*{The two-weak case and the empty cells}

Suppose first that \(C,\alpha\) are weak.  The actual proof
Section~\ref{component:p534f57e1da6f} supplies the following source implications,
without a tensor-SOS premise.  The coefficient-field experiments
which kill three matching packets show that
\[
                    n,m,A,D\text{ are unbounded}.             \tag{CT7}
\]
For example the remaining \(D\) experiment has full value
\(4D\varepsilon^2-2\varepsilon+4(2+\varepsilon)C
+8(1+\varepsilon)\alpha\).
If \(D\) were bounded, a sufficiently small positive real
\(\varepsilon\) would contradict physicality.  The other three
experiments are its exact event-cycle copies.
The exact signed-core estimate
\[
 H\geq-2u_cv_c\|\psi_c\|^2
\]
follows from the printed matrix-square identity
\((M_H-LI)^2=L^2I+4u_cv_cLJ_0\), \(J_0^2=I\).
Each matching amplitude divided by \(\sqrt{u_cv_c}\) is therefore
infinitesimal at a norm-one contact.  Compact angle coordinates
and, if needed, CT4 imply
\[
 c\text{ bounded},\qquad a,b,t\sim1,\qquad
 r,h,g,j\text{ positive infinitesimal}.                    \tag{CT8}
\]
In the small-angle case the two matching equations divided by
\(\sqrt{u_cv_c}\) first force \(v_c/u_c\sim1\) and \(W_c/Z_c\sim1\);
this is a relative bound, not just a compact limiting statement.
The same body proves \(\alpha/u_c^2\) infinitesimal before using
scaled stationarity.  Its invertible four-amplitude Jacobian has
residue determinant \(2\), and the weighted stationary gradients
give \(n\ell_n,A\ell_A,m\ell_m,D\ell_D\) residues
\((-1,1,1,-1)/4\).  Thus all strict signs in CT8 are justified.

For \(c\) infinitesimal use the complete six-root radial construction
in Section~\ref{component:p9f1235511111}, Section~\ref{component:p53752031d16c}, and
Section~\ref{component:p163e0ebf59f6}.  Only its radial domain CT8 is imported.
Its pure-source weight \(N\) satisfies
\[
 N\asymp r+h+g+j,\qquad
 r/N,h/N,g/N,j/N=O(1),\qquad
 \operatorname{res}\bigl(4\alpha/(c^4N)\bigr)=1.          \tag{CT9}
\]
The corrected radial root is \(\sqrt N(u^2-c^2)\).
The actual proof retains all four normalized event roots and every
state normal; the final exact current-zero correction makes every
radial cross coefficient belong to \((r,h,g,j)\) before division
by \(\sqrt N\).  Its Gram has strict residue on the full six-root
list, including the complementary \(W\)-state chart.
The actual outer proof Section~\ref{component:pc7dacaa3258e}
uses
\[
             1/(Cn),\ 1/(Cm),\ 1/(CA),\ 1/(CD)=O(1),
\]
which follow from
CT9 and the unchanged source formulas.  Its reciprocal-corner
conductor and its two complete state charts retain the weak \(C\)
packet exactly.
For \(c\) a unit, Section~\ref{component:p8b53b68d8af6} gives
\(\alpha/C\asymp1\), both sets of four product bounds, and the
same outer proof after CT4 covers the remaining corner.
Both cases use the detailed arbitrary-scale physical-source
noncommuting-residue theorem A6 for the noncommuting testing
complement.  A6's full geometry, curvature, weighted matching,
promotion, and induction bodies are explicit dependencies.
This completes every two-weak source, rather than assuming a
preselected radial chart.

For the right column of CT6, the exact proof
Section~\ref{component:pe56fb394b522} establishes the II, IU, UI and UU
contradictions from CT2 and \(n\geq0\).
Some useful details fix its scope.  On II, putting
\[
 H=c^2h,\quad L=c^2(b-h),\quad B=H+L
\]
in its own actual-angle notation makes the positive numerator
bound force \(\alpha\) weak whenever \(C\) is weak.  On UI, the
same source numerator and the \(n\geq0\) inequality force
\(r\) infinitesimal and \(h\) a unit, contradicting weakness
of \(C\).  On IU, bounded \(b\) is a unit by the inequality in
CT10 below and gives \(C\) strong; unbounded \(b\) gives
\(C\asymp L^{-1}\) with \(L\) bounded.  On UU, \(a,c\) are units,
\(\Delta\) and \(b+r\) are units, and the positive numerator
identity makes \(C,\alpha\) weak together or strong together.
These are necessary-source arguments, not all-tests conclusions.

The OI and OU contradictions use
Section~\ref{component:pb7189b9c1f6e}.
With \(q=u_c^{-2}\) infinitesimal and \(y=v_c^2\) bounded, its
exact numerator for \(n\) is \(-qP/(2rQ)\), \(Q>0\).
The printed positive remainder gives
\[
 P\geq r^2(1+y)(H+yr)-Tqr(y+2)-q^2(T^2+T),\qquad
 0<T<H+y .
\]
If \(q/r\) were infinitesimal, division by \(r^2(H+yr)\)
would force \(P>0\), contradicting \(n\geq0\).
Thus \(r/q\) is bounded, \(H\sim1\), and the exact cancelled
\(C\) formula has a positive-unit numerator and bounded
positive denominator.  Consequently \(C\) is strong.  Literal
\(n=0\) obeys the same contradiction.

\subsubsection*{The single-angle join has no omitted opening}

For an IU source put \(x=u_c^2\) infinitesimal and
\(k=v_c\) a unit.  If \(b\) is bounded, the exact source
inequality \(n\geq0\) gives
\[
 r_0\{b_0+k_0^2(b_0+1)\}\leq b_0,\qquad
 r_0\geq1-b_0,\qquad
 b_0\geq k_0/\sqrt{1+k_0^2}>0 .                       \tag{CT10}
\]
The actual body Section~\ref{component:p7294aae40ee2} proves
the last bound also when \(b_0\geq1\).
There are then exactly three source alternatives:
\[
\begin{array}{c|l}
g\text{ unit}&\text{three-opening proof, arbitrary }r,h,b-h>0,\\
g\text{ infinitesimal},\ r\text{ unit}
             &\text{last-opening proof},\\
g,r\text{ infinitesimal}
             &\text{joint centered-packet proof}.
\end{array}\tag{CT11}
\]
All bounded positive \(g,r\) have been included.
The first proof supplies both bounded and unbounded \(n\)
normalizations; the second uses the triangular hard normal of
the entire \(m\) amplitude; the third retains both contact-zero
centered \(n,m\) packets.  Their common full scalar proof is the
six-term balanced Laurent decomposition in
Section~\ref{component:pec59340d644f}, with its five-term limiting
version and joint version proved in the respective bodies.
This supplies global equality uniqueness and a positive contact
Hessian.  The same bodies print complete first-normal divisions,
hard blocks, the literal \(n=0\) double-zero corner and all four
outer-core identities.  Their remaining testing complement uses
the actual detailed A6 theorem.  No normal is omitted and no
bound on \(h\) or \(b-h\) is restored implicitly.

If \(b\) is unbounded, put \(B=xb,H=xh,L=x(b-h)\).
The original gaps give
\[
 B=H+L,\quad 0<H<1,\quad 0<L<k^2(1+xr).
\]
Thus \(B\) is bounded.  The exact proof
Section~\ref{component:pb60046a37191} cancels the original
denominators and gives
\[
 1-r\asymp1,\quad m\asymp b^{-2},\quad
 C\asymp L^{-1},\quad A\asymp H^{-1},\quad
 D\asymp B/j,\quad \alpha\asymp B/(1-H).               \tag{CT12}
\]
If \(B\) is a unit, \(C,\alpha\) are strong and the full strong
pitchfork below applies.  If \(B\) is infinitesimal, the actual
Section~\ref{component:pab6222b899ef} theorem applies with
\(\sigma=1/b\), with no restriction on \(B/\sigma\).
Its separate homogeneous scalar lemma proves the entire
contact-ideal lift on
\[
 (uW-(B/k)Z)/\sqrt B,\quad uZ-\sigma kW,\quad
 \sqrt B(v-k)Z,\quad u(v-k)Z .
\tag{CT13}
\]
The whole-parent proof retains the possibly nonzero first-normal
angular coupling and buys its strict pivot by the exact zero
\[
 T^2-\sigma U^2-ST(v-k)-\sigma kRU(v-k)=0 .
\]
It proves both full normals, arbitrary \(W/Z\), every \(n\) scale,
all openings, and the literal \(n=0\) corner.  Its outer charts
use the displayed general-single-angle identities with scale
\(\sqrt B\), and the detailed A6 theorem for the final complement.
This exhausts IU.  CT4 exhausts UO with no new source assumption.

The strong pitchfork invoked in CT12 is itself complete at all
upper scales.  Its three-large-event body covers bounded-unit
\(C,\alpha\), with \(n,A,D\) arbitrary; its two-large-event
predecessor supplies the printed stationary equations, full
normal Taylor division and real boundary Schur calculation.
The body Section~\ref{component:pf7943b34a7dc} treats
unbounded \(C\) or \(\alpha\).  Its coefficient-field state
equations force \(m/k^4\asymp1\) and bounded weighted graph
coordinates before any small gap is inverted.  It uses the
actual large \(n,C,\alpha\) or \(n,D,\alpha\) packets as normals.
The five-root full-parent Gram has a strict angular Schur pivot
because the normalized real boundary numerator is
\(q_{12}(v^2-a_0^2)^2\) with \(q_{12}>0\).
Only this real residue Schur computation uses a determinant
inverse; no mixed tensor denominator is inverted.
For its outer charts one may use the complete cap and corner
proof printed in that same strong-complement body, replacing
the older predecessor's reference to a preceding outer corollary.

\subsubsection*{The remaining occupied cells}

\paragraph{II and OO with both complementary events strong.}
The exact common-II source reduction proves that
\(\alpha\) is strong precisely when \(B=H+L\) is a unit.
Then \(L\) is infinitesimal, \(H\) and \(1-r\) are units, and
\[
 m/z^2,\ A/y,\ CL,\ Dyj,\ \alpha(1-H)\asymp1,
 \qquad nr=O(1).                                      \tag{CT14}
\]
The exact numerator at \(r=1\) is negative, giving the asserted
unit gap \(1-r\); none of \(r,j,1-H\) is assumed a unit.
The proof Section~\ref{component:pe16499ccc4c6} retains all four centered
large-packet excesses and gives the full three-normal remainder.
Its angular identity, harmonic \(C,D\) square and free radial
pivot are exact; the normalized hard residue and the remaining
two-by-two Schur margin are strictly positive.
The separate Section~\ref{component:p590112f102bd} uses the original
homogeneous parent and harmonic current, not a higher-degree
homogenization of the local graph.  It proves all remaining
angles and state charts, including the endpoint forms.
Thus there is no local-to-remote scope gap.  CT4 gives OO.

\paragraph{IO, weak or strong \(C\).}
The weak IH source entry proves \(B\) unbounded and
\(\rho=B/y\) bounded with \(0\leq\bar\rho\leq1\).
The full interior-corner proof, both \(\rho=0,1\) source-endpoint
extensions, remote-chart proof, and separate open-axis proof
cover all independent angles and states.  The latter open-axis
proof is indispensable: it supplies the product-pivot split
when \(u\) is a unit and \(1/v\) is infinitesimal.
All these bodies retain \(nr\) and the full \(D\) packet at the
closing endpoint; literal \(n=0\) is included.
The actual weak-IH all-tests join merely assembles these
already specified bodies.

If \(C,\alpha\) are both strong, the separate strong-angle
source entry gives \(B\) a bounded unit, \(r=O(1/y)\),
\(Dy^2\asymp1\) and the remaining opening-safe products.
The strong-IH local proof reuses only the explicit polynomial
current identities of the weak-IH proof, with its own source
estimates and finite \(C\) cap; it does not import the weak-IH
source conclusion.  Its remote proof includes the two open
axes and reciprocal corner and uses detailed A6 only when
both testing pairs are noncommuting in residue.

\paragraph{UI and OU with both complementary events strong.}
The actual Section~\ref{component:pc3daa35e9a9b} in the C2-UI
directory cancels all closing source openings before dividing
the normals.  In particular \(nr,m\mu,\alpha\pi,C\rho k^2,
D\sigma k^2\), and the difference \(Ct-Dj\), are bounded.
The exact state inverse is polynomial and never divides by
a state coordinate.  The reciprocal-contact lemma and local
certificate prove the full contact Gram, including all three
normals.  The complete Section~\ref{component:p6b22e944bd7c}
proof supplies real interior-source approximation for its
finite-cap step and strict Grams on every complementary
axis and state chart.  Its Laurent equality and positive
contact Hessian follow from the actual source equations.
No new global source classification is imported from this
local theorem.  CT4 gives OU.

\paragraph{OI.}
The exact HI source entry proves directly from \(n\geq0\)
that \(R=r/q\) is bounded, \(H\sim1\), and
\[
 mqK,\ Cyt,\ Dyj,\ \alpha qE,\ A/(qy)\asymp1,
 \quad q=u_c^{-2},\ y=v_c^2\text{ infinitesimal}.       \tag{CT15}
\]
Thus \(m,C,D,\alpha\) are all unbounded, without a prior
strong-event assumption.  Its remote full-parent theorem
caps these four coefficients at a sufficiently large
\emph{real} constant after fixing the testing residue.
It prints strict full state forms on every remote chart
and literal critical axis; only \(1/u,v\) both nonzero
infinitesimal remain.
The retained-contact module gives the complete local
four-monomial circuit and radical conductor there.
The full HI module includes the actual \(n\) packet, its
two complete normals, and the endpoint \(R_0=0\) where
\(n\) can be unbounded.  Its constant circuit multiplier
lies in \([2/3,2]\), and its normalized hard residue
is strictly positive.  Its proof uses the predecessor
only for the explicit half-angle conductor, never for
an unsupported \(n=0\)-to-\(n>0\) inference.
Together these three bodies close OI.

\paragraph{UU and the five-event terminal.}
The exact source-angle equations give \(m,A,D\) strong on
the both-strong UU cell.  The actual five-retained-events
proof uses \(\mu=\min(1,m,C,A,D,\alpha)\), a positive unit,
and retains all excess original packets.
Its commuting-axis full state forms and three nontrivial
double-corner coordinate Grams are strict.  For
noncommuting testing pairs it uses the detailed arbitrary-scale
physical-source theorem A6 for the original source.
The selected capped parent is never asserted globally
physical.  This distinction prevents a contact-zero
packet from being improperly subtracted at the source
contact.

\paragraph{OO in the directed weak case.}
CT4 gives exactly the native II domain HJ1:
\(C\) strong, \(\alpha\) weak, both actual angles
infinitesimal, and \(\delta_{\rm tot}\) infinitesimal.
The current directed proof in
Section~\ref{component:pec6c617d3a3d} derives HJ1--HJ16 from the
unchanged source equations. It proves the equivalence of the weak
face with \(B\) infinitesimal, the source-only earlier exits, and
the exhaustive ten-row partition with each leaf's actual hypotheses.

Its R1 and R2 rows use the complete current R1N and R2 joins;
R3/R4 split into unit-\(e\) and EN cases; R5--R8 use the displayed
bounded/unbounded source splits and complete local and remote bodies.
The regular unit-\(e\) exits are WA, WE and SL. Every remaining
bounded-\(a\), small-\(e\) source, including the former WC and R9
cases, is covered by BA1--BA15 in
Section~\ref{sec:new-bounded-a-weak-c2}.
For unbounded \(a\), the R10 alternatives are
\(n/e^2\) noninfinitesimal, \(0<n/e^2\) infinitesimal, and literal
\(n=0\). Their proofs are respectively the weighted full-parent
identity, the positive-thin contact proof with its unit-\(W\) state
complement, and the exact boundary map OP4 followed by the retained
native C1 small-\(e\) or centered-\(A\) body.

The HJ derivation verifies these source conditions before any leaf
is used. The leaves retain the complete parent, all state normals,
and their projective testing joins. The \(n=0\) C1 uses are the two
specified local all-tests bodies under their exact native map;
no negative formal-beta source is declared physical. Consequently
this cell uses neither an omitted bare-source audit nor this global
C2 theorem, the whole C1 theorem, or AU as a premise.

\subsubsection*{Proof-use scope and the conclusion}

Every mathematical input used above is the explicit proof body at
its displayed source hypotheses, including the single-angle and
homogeneous scalar arguments, the complete normal and source-endpoint
constructions, and the noncommuting residue proofs. The retained
formulas and their bounds provide the following directed assembly:
\[
\begin{gathered}
 \text{native source entry: EN;}\\
 \text{positive total scale: AU;}\\
 \text{directed weak II: HJ with BA and the explicit R10 split.}
\end{gathered}\tag{CT16}
\]
Only these actual source hypotheses and complete full-parent
constructions are used below.
The bodies in CT10--CT15 use the exact source equations,
positivity, the typed source scales and the displayed
opening bounds.  Their scalar divisions, normal bounds,
caps and outer joins do not use exclusion from any
blanket list of earlier admissions.
Those exclusions are therefore absent from the theorem:
this is a proof-use verification on the actual bodies,
not deletion of a logical hypothesis by relabelling.

Every positive \(\delta_{\rm tot}\) has one of the two
types in CT3.  In the infinitesimal case \(C,\alpha\)
have exactly one of the three weak/strong patterns
used above; CT4 orients the asymmetric pattern.
The nine ordered angle types in CT6 are exhaustive.
Every occupied cell has a full-parent theorem at its
exact source hypotheses, and every empty cell has the
necessary-source contradiction stated above.
Thus CT1 is tensor SOS at every independent test.
Complex states split into real and imaginary parts
because CT1 is a real quadratic form in the state.
All source and testing normalizations are pure-field
square factors.  Restoring the original positive \(d\)
completes the theorem.

